% Generated draft: requires mathematical and source review.
\documentclass[11pt]{article}
\usepackage[margin=25mm]{geometry}
\usepackage{fontspec}
\setmainfont{Latin Modern Roman}
\usepackage{amsmath,amssymb,mathtools,unicode-math}
\setmathfont{Latin Modern Math}
\usepackage{enumitem,needspace,xurl,hyperref}
\setlength{\parindent}{0pt}
\setlength{\parskip}{5pt}
\setlength{\emergencystretch}{3em}
\newcommand{\N}{\mathbb{N}}
\newcommand{\Z}{\mathbb{Z}}
\newcommand{\Q}{\mathbb{Q}}
\newcommand{\R}{\mathbb{R}}
\newcommand{\C}{\mathbb{C}}
\newcommand{\F}{\mathbb{F}}
\newcommand{\A}{\mathbb{A}}
\newcommand{\PP}{\mathbb{P}}
\newcommand{\E}{\mathbb{E}}
\newcommand{\eps}{\varepsilon}
\newcommand{\dd}{\,\mathrm{d}}
\begin{document}
\section*{UnsolvedMath ID: 30003051}
\textbf{Source problem code:} catalogue entry
\subsection{Definitions and mathematical statement}
The supplied record proposes a state-dependent map for a system with matrix set \(\mathcal M=\{A_0,A_1\}\). It does not specify the matrix dimension, scalar field, definition of \(\rho(\mathcal M)\), or convention for \(\widehat{\mathcal M}^{*}\). The following makes a conventional interpretation explicit, conditionally; these choices require source verification.

Let \(d\geq 1\), let \(\mathbb K\in\{\R,\C\}\), and suppose \(A_0,A_1\in\mathbb K^{d\times d}\), acting on \(X=\mathbb K^d\) with Euclidean norm \(\lVert\cdot\rVert_2\). If \(\rho(\mathcal M)\) denotes the joint spectral radius, its conventional definition is
\[
\rho(\mathcal M)=\limsup_{n\to\infty}\left(\max_{B_1,\ldots,B_n\in\mathcal M}\lVert B_n\cdots B_1\rVert_2\right)^{1/n},
\]
where the matrix norm is the induced operator norm. The displayed normalization requires \(\rho(\mathcal M)>0\), a condition not explicitly stated in the supplied formulation. Subject to that requirement, put
\[
\widehat{\mathcal M}=\left\{\frac{A_0}{\rho(\mathcal M)},\frac{A_1}{\rho(\mathcal M)}\right\}.
\]
Under the convention that the star denotes all finite products, including the empty product, write
\[
\widehat{\mathcal M}^{*}=\{I_d\}\cup\bigcup_{n\geq1}\{B_n\cdots B_1:B_1,\ldots,B_n\in\widehat{\mathcal M}\}.
\]
The inclusion of the empty product is a convention requiring confirmation. Retain the source hypothesis that \(\widehat{\mathcal M}^{*}\) is bounded; in this interpretation this means
\[
\sup_{B\in\widehat{\mathcal M}^{*}}\lVert B\rVert_2<\infty.
\]
Define
\[
p(x)=\sup_{B\in\widehat{\mathcal M}^{*}}\lVert Bx\rVert_2,
\qquad
f(x)=\begin{cases}
A_0x,&p(x)=p(A_0x),\\
A_1x,&p(x)\ne p(A_0x).
\end{cases}
\]
This retains the source's use of the original matrices \(A_0,A_1\) in the branches and of \(A_0\), without normalization, in the equality test. The boundedness hypothesis concerns products of the normalized matrices.

For an algebraic interpretation, let \(\mathcal F(X)=\{g:X\to\C\}\). The composition operator induced by \(f\) is
\[
U_f:\mathcal F(X)\to\mathcal F(X),\qquad (U_fg)(x)=g(f(x)).
\]
This supplies a definition on all scalar observables only. The record does not specify the intended invariant observable space, topology, norm, or measure, so it does not determine a particular analytic realization of the Koopman operator.

The source asks whether the Koopman operator induced by this state-dependent map is an appropriate operator for the switching system and what useful information about that system it provides. Neither the admissible switching signals nor a mathematical criterion for appropriateness, usefulness, or information preservation is specified. Thus the supplied target is an open modeling question, rather than a fully quantified conjecture. The conditional definitions above make the proposed construction explicit but do not resolve those missing specifications or assert an affirmative answer.
\subsection{Short English statement}
The proposal chooses between two matrix updates using an equality of suprema over normalized matrix products, then studies observables by composition with the resulting map. The question asks whether this Koopman construction appropriately describes the switching system and yields useful information. The state space, normalization conventions, observable space, admissible switching rules, and criterion of success require clarification.
\subsection{Sources}
\begin{itemize}
\item Local UnsolvedMath registry, supplied complete record with permanent ID \(30003051\), titled \emph{State-Dependent Koopman Operators for Switching Systems}; exact source problem code \texttt{catalogue entry}. The record belongs to the Oberwolfach Reports Open Problems collection and lists proposed year 2016; no individual proposer is supplied.
\item Supplied source citation: \emph{Mini-Workshop: Applied Koopmanism} (2016), Oberwolfach Reports, pp.\ 297--340, DOI \texttt{10.4171/owr/2016/7}. \url{https://doi.org/10.4171/owr/2016/7}. The background describes these reports as documentation of mathematical workshops and their open-problem discussions.
\item Supplied literature source labeled ``original OWR report'': \url{https://doi.org/10.4171/owr/2016/7}.
\item Supplied literature source labeled ``2025 switched-systems Koopman paper'': \url{https://doi.org/10.1016/j.sysconle.2025.106340}. No authors, title, or further bibliographic details are supplied for this item.
\item Literature assessment dated exactly \texttt{2026-08-21}: the supplied curation labels the problem open and describes the appropriateness and information-preserving character of the proposed operator as a modeling question; it states that current applications do not establish a general affirmative theorem. This is supplied dated curation, not independently verified here. The background expressly requests independent verification of time-sensitive status and cited claims.
\item The registry labels the statement \texttt{corrected\_verified} and reports that it was checked against the official Oberwolfach report, preserving its hypotheses, notation, and mathematical target. These are supplied provenance claims, not an independent verification performed for this draft.
\end{itemize}
\end{document}
